% Part: first-order-logic
% Chapter: tableaux
% Section: proving-things

\documentclass[../../../include/open-logic-section]{subfiles}

\begin{document}

\iftag{FOL}
      {\olfileid{fol}{tab}{pro}}
      {\olfileid{pl}{tab}{pro}}

\olsection{\usetoken{P}{tableau}નાં ઉદાહરણો}

\begin{ex}
$(!A \land !B) \lif !A$ !!{sentence} માટે બંધ !!{tableau} શોધીએ.

આપણે !!{tableau}ની ટોચે તેને અનુરૂપ ધારણા લખીને શરૂઆત કરીએ છીએ.
\begin{oltableau}
  [\sFmla{\False}{(\formula{A} \land \formula{B}) \lif \formula{A}},
    just = \TAss]
\end{oltableau}

માત્ર એક ધારણા છે, તેથી નિયમ લાગુ કરી શકાય એવું માત્ર એક
!!{signed formula} છે. (દરેક !!{signed formula} માટે લાગુ કરી શકાય એવો
વધુમાં વધુ એક જ નિયમ હોય છે: તે !!{sentence}ના ચિહ્ન અને
!!{main operator}ને અનુરૂપ નિયમ છે.) આ કિસ્સામાં આપણે
$\TRule{\False}{\lif}$ લાગુ કરવો પડે છે.
\begin{oltableau}
  [\sFmla{\False}{(\formula{A} \land \formula{B}) \lif \formula{A}},
    checked, just = \TAss
    [\sFmla{\True}{\formula{A} \land \formula{B}},
      just={\TRule{\False}{\lif}[1]}
      [\sFmla{\False}{\formula{A}}, just={\TRule{\False}{\lif}[1]}]
    ]
  ]
\end{oltableau}
કયા !!{signed formula}sને તેમના અનુરૂપ નિયમો લાગુ કર્યા છે તેનો હિસાબ
રાખવા માટે આપણે વાક્યની બાજુમાં ખરાનું ચિહ્ન મૂકીએ છીએ. પરંતુ નિયમ
બધી ખુલ્લી શાખાઓ પર લાગુ કર્યો હોય \emph{તો જ} ખરાનું ચિહ્ન મૂકો. કોઈ
!!a{signed formula}ને દરેક ખુલ્લી શાખામાં અનુરૂપ નિયમ લાગુ થઈ ગયા પછી
આપણે તેની પાસે પાછા ફરીને નિયમ ફરી લાગુ કરવો પડતો નથી. આ કિસ્સામાં
માત્ર એક જ શાખા છે, તેથી નિયમ પણ માત્ર એક વાર લાગુ કરવો પડે છે. (નોંધો
કે ખરાનાં ચિહ્નો ટેબ્લો રચવાની સગવડ માત્ર છે અને ટેબ્લોના વાક્યવિન્યાસનો
વિધિવત્ ભાગ નથી.)

હવે નિયમ લાગુ કરી શકાય એવું એક નવું !!{signed formula} છે: પંક્તિ~$2$
પરનું $\sFmla{\True}{!A \land !B}$. $\TRule{\True}{\land}$ નિયમ લાગુ
કરતાં આ મળે છે:
\begin{oltableau}
  [\sFmla{\False}{(\formula{A} \land \formula{B}) \lif \formula{A}},
    checked, just = \TAss
    [\sFmla{\True}{\formula{A} \land \formula{B}},
      just={\TRule{\False}{\lif}[1]}, checked
      [\sFmla{\False}{\formula{A}}, just={\TRule{\False}{\lif}[1]}
        [\sFmla{\True}{\formula{A}}, just={\TRule{\True}{\land}[2]}
          [\sFmla{\True}{\formula{B}}, just={\TRule{\True}{\land}[2]}, close
          ]
        ]
      ]
    ]
  ]
\end{oltableau}
હવે શાખામાં $\sFmla{\True}{!A}$ (પંક્તિ~$4$ પર) અને
$\sFmla{\False}{!A}$ (પંક્તિ~$3$ પર) બંને હોવાથી શાખા બંધ છે. આ એકમાત્ર
શાખા હોવાથી !!{tableau} બંધ છે. આમ, આપણે~$(!A \land !B) \lif !A$ માટે
બંધ !!{tableau} શોધી લીધો છે.
\end{ex}

\begin{ex}
હવે $(\lnot !A \lor !B) \lif (!A \lif !B)$ માટે બંધ !!{tableau}
શોધીએ.

આપણે તેને અનુરૂપ ધારણાથી શરૂઆત કરીએ છીએ:
\begin{oltableau}
  [\sFmla{\False}{(\lnot \formula{A} \lor \formula{B}) \lif
      (\formula{A} \lif \formula{B})}, just=\TAss]
\end{oltableau}
આ !!{tableau}ના એકમાત્ર !!{signed formula}નો !!{main operator}~$\lif$
અને ચિહ્ન~$\False$ છે, તેથી તેને $\TRule{\False}{\lif}$ નિયમ લાગુ કરતાં
આ મળે છે:
\begin{oltableau}
  [\sFmla{\False}{(\lnot \formula{A} \lor \formula{B})
      \lif (\formula{A} \lif \formula{B})}, just=\TAss, checked
    [\sFmla{\True}{\lnot \formula{A} \lor \formula{B}},
      just={\TRule{\False}{\lif}[1]}
      [\sFmla{\False}{(\formula{A} \lif \formula{B})},
        just={\TRule{\False}{\lif}[1]}
      ]
    ]
  ]
\end{oltableau}
હવે પંક્તિ~$2$ને~$\TRule{\True}{\lor}$ લાગુ કરવો કે પંક્તિ~$3$ને
$\TRule{\False}{\lif}$ લાગુ કરવો તે પસંદ કરી શકીએ છીએ. દરેક
!!{signed formula}ને દરેક શાખામાં તેનો અનુરૂપ નિયમ લાગુ કરીએ ત્યાં સુધી
આપણે કયો ક્રમ પસંદ કરીએ તે વાસ્તવમાં મહત્ત્વનું નથી. તેથી પહેલો વિકલ્પ
પસંદ કરીએ. $\TRule{\True}{\lor}$ નિયમ !!{tableau}ની શાખાઓ પાડવાની છૂટ
આપે છે અને નિયમનાં બે ફલિતો નવી બે શાખાઓમાં ઉમેરાયેલાં નવાં
!!{signed formula}s બનશે. તેનાથી આ મળે છે:
\begin{oltableau}
  [\sFmla{\False}{(\lnot \formula{A} \lor \formula{B}) \lif
      (\formula{A} \lif \formula{B})}, just=\TAss, checked
    [\sFmla{\True}{\lnot \formula{A} \lor \formula{B}},
      just={\TRule{\False}{\lif}[1]}, checked
      [\sFmla{\False}{(\formula{A} \lif \formula{B})},
        just={\TRule{\False}{\lif}[1]}
        [\sFmla{\True}{\lnot \formula{A}}, just={\TRule{\True}{\lor}[2]}]
        [\sFmla{\True}{\formula{B}}, just={\TRule{\True}{\lor}[2]}]
      ]
    ]
  ]
\end{oltableau}
આપણે પંક્તિ~$3$ને હજી $\TRule{\False}{\lif}$ નિયમ લાગુ કર્યો નથી; હવે
તેમ કરીએ. સમય બચાવવા માટે તેને બંને શાખાઓ પર લાગુ કરીએ છીએ. યાદ કરો કે
દરેક ખુલ્લી શાખામાં અનુરૂપ નિયમ લાગુ કર્યો હોય તો જ આપણે
!!a{signed formula}ની બાજુમાં ખરાનું ચિહ્ન મૂકીએ છીએ. તેથી નિયમ જે
!!{signed formula}ને લાગુ પડે તે ધરાવતી દરેક શાખાના છેડે નિયમ લાગુ કરવો
સારો ઉપાય છે. એમ કરવાથી જુદી જુદી શાખાઓમાં નીચે ગયા પછી એ જ
!!{signed formula} પાસે પાછા ફરવું નહીં પડે.
\begin{oltableau}
  [\sFmla{\False}{(\lnot \formula{A} \lor \formula{B}) \lif
      (\formula{A} \lif \formula{B})}, just=\TAss, checked
    [\sFmla{\True}{\lnot \formula{A} \lor \formula{B}},
      just={\TRule{\False}{\lif}[1]}, checked
      [\sFmla{\False}{(\formula{A} \lif \formula{B})},
        just={\TRule{\False}{\lif}[1]}, checked
        [\sFmla{\True}{\lnot \formula{A}}, just={\TRule{\True}{\lor}[2]}
          [\sFmla{\True}{\formula{A}}, just={\TRule{\False}{\lif}[3]}
            [\sFmla{\False}{\formula{B}}, just={\TRule{\False}{\lif}[3]}]
          ]
        ]
        [\sFmla{\True}{\formula{B}}, just={\TRule{\True}{\lor}[2]}
          [\sFmla{\True}{\formula{A}}, just={\TRule{\False}{\lif}[3]}
            [\sFmla{\False}{\formula{B}}, just={\TRule{\False}{\lif}[3]}, close]
          ]
        ]
      ]
    ]
  ]
\end{oltableau}
હવે જમણી શાખા બંધ છે. ડાબી શાખામાં પંક્તિ~$4$ને હજી
$\TRule{\True}{\lnot}$ નિયમ લાગુ કરી શકીએ છીએ. તેનાથી
~$\sFmla{\False}{!A}$ મળે છે અને ડાબી શાખા પણ બંધ થાય છે:
\begin{oltableau}
  [\sFmla{\False}{(\lnot \formula{A} \lor \formula{B}) \lif
      (\formula{A} \lif \formula{B})}, just=\TAss, checked
    [\sFmla{\True}{\lnot \formula{A} \lor \formula{B}},
      just={\TRule{\False}{\lif}[1]}, checked
      [\sFmla{\False}{(\formula{A} \lif \formula{B})},
        just={\TRule{\False}{\lif}[1]}, checked
        [\sFmla{\True}{\lnot \formula{A}}, just={\TRule{\True}{\lor}[2]}
          [\sFmla{\True}{\formula{A}}, just={\TRule{\False}{\lif}[3]}
            [\sFmla{\False}{\formula{B}}, just={\TRule{\False}{\lif}[3]}
              [\sFmla{\False}{\formula{A}},
                just={\TRule{\True}{\lnot}[4]}, close]
            ]
          ]
        ]
        [\sFmla{\True}{\formula{B}}, just={\TRule{\True}{\lor}[2]}
          [\sFmla{\True}{\formula{A}}, just={\TRule{\False}{\lif}[3]}
            [\sFmla{\False}{\formula{B}},
              just={\TRule{\False}{\lif}[3]}, close]
          ]
        ]
      ]
    ]
  ]
\end{oltableau}
\end{ex}

\begin{ex}
ધારણાઓ તરીકે ગમે તેટલાં !!{signed formula}s માટે આપણે !!{tableau}s આપી
શકીએ છીએ. ઘણી વાર શાખાઓ પાડતા એક કરતાં વધુ નિયમો લાગુ કરવા પણ જરૂરી
હોય છે; અને સામાન્ય રીતે !!a{tableau}ને ગમે તેટલી શાખાઓ હોઈ શકે છે.
દાખલા તરીકે, આ ગણ માટેનો !!a{tableau} લો:
$\{\sFmla{\True}{!A \lor (!B \land !C)}, \sFmla{\False}{(!A \lor !B) \land (!A
  \lor !C)}\}$. આપણે પહેલી ધારણાને $\TRule{\True}{\lor}$ લાગુ કરીને
શરૂઆત કરીએ છીએ:
\begin{oltableau}
  [\sFmla{\True}{\formula{A} \lor (\formula{B} \land \formula{C})},
    just=\TAss, checked
    [\sFmla{\False}{(\formula{A} \lor \formula{B}) \land
        (\formula{A} \lor \formula{C})}, just=\TAss
      [\sFmla{\True}{\formula{A}}, just={\TRule{\True}{\lor}[1]}]
      [\sFmla{\True}{\formula{B} \land \formula{C}},
        just={\TRule{\True}{\lor}[1]}]
    ]
  ]
\end{oltableau}
હવે પંક્તિ~$2$ને $\TRule{\False}{\land}$ નિયમ લાગુ કરી શકીએ છીએ. આપણે
બંને શાખાઓ પર એકસાથે તેમ કરીએ છીએ અને તેથી પંક્તિ~$2$ પર ખરાનું ચિહ્ન
મૂકી શકીએ છીએ:
\begin{oltableau}
  [\sFmla{\True}{\formula{A} \lor (\formula{B} \land \formula{C})},
    just=\TAss, checked
    [\sFmla{\False}{(\formula{A} \lor \formula{B}) \land
        (\formula{A} \lor \formula{C})}, just=\TAss, checked
      [\sFmla{\True}{\formula{A}}, just={\TRule{\True}{\lor}[1]}
        [\sFmla{\False}{\formula{A} \lor \formula{B}},
          just={\TRule{\False}{\land}[2]}]
        [\sFmla{\False}{\formula{A} \lor \formula{C}},
          just={\TRule{\False}{\land}[2]}]
      ]
      [\sFmla{\True}{\formula{B} \land \formula{C}},
        just={\TRule{\True}{\lor}[1]}
        [\sFmla{\False}{\formula{A} \lor \formula{B}},
          just={\TRule{\False}{\land}[2]}]
        [\sFmla{\False}{\formula{A} \lor \formula{C}},
          just={\TRule{\False}{\land}[2]}]
      ]
    ]
  ]
\end{oltableau}
હવે $!A \lor !B$ ધરાવતી બધી શાખાઓને $\TRule{\False}{\lor}$ લાગુ કરી
શકીએ છીએ:
\begin{oltableau}
  [\sFmla{\True}{\formula{A} \lor (\formula{B} \land \formula{C})},
    just=\TAss, checked
    [\sFmla{\False}{(\formula{A} \lor \formula{B}) \land
        (\formula{A} \lor \formula{C})}, just=\TAss, checked
      [\sFmla{\True}{\formula{A}}, just={\TRule{\True}{\lor}[1]}
        [\sFmla{\False}{\formula{A} \lor \formula{B}},
          just={\TRule{\False}{\land}[2]}, checked
          [\sFmla{\False}{\formula{A}}, just={\TRule{\False}{\lor}[4]}
            [\sFmla{\False}{\formula{B}},
              just={\TRule{\False}{\lor}[4]}, close]
          ]
        ]
        [\sFmla{\False}{\formula{A} \lor \formula{C}},
          just={\TRule{\False}{\land}[2]}]
      ]
      [\sFmla{\True}{\formula{B} \land \formula{C}},
        just={\TRule{\True}{\lor}[1]}
        [\sFmla{\False}{\formula{A} \lor \formula{B}},
          just={\TRule{\False}{\land}[2]}, checked
          [\sFmla{\False}{\formula{A}}, just={\TRule{\False}{\lor}[4]}
            [\sFmla{\False}{\formula{B}},
              just={\TRule{\False}{\lor}[4]}]
          ]
        ]
        [\sFmla{\False}{\formula{A} \lor \formula{C}},
          just={\TRule{\False}{\land}[2]}]
      ]
    ]
  ]
\end{oltableau}
હવે સૌથી ડાબી શાખા બંધ છે. હવે $!A \lor !C$ને
$\TRule{\False}{\lor}$ લાગુ કરીએ:
\begin{oltableau}
  [\sFmla{\True}{\formula{A} \lor (\formula{B} \land \formula{C})},
    just=\TAss, checked
    [\sFmla{\False}{(\formula{A} \lor \formula{B}) \land
        (\formula{A} \lor \formula{C})}, just=\TAss, checked
      [\sFmla{\True}{\formula{A}}, just={\TRule{\True}{\lor}[1]}
        [\sFmla{\False}{\formula{A} \lor \formula{B}},
          just={\TRule{\False}{\land}[2]}, checked
          [\sFmla{\False}{\formula{A}},
            just={\TRule{\False}{\lor}[4]}
            [\sFmla{\False}{\formula{B}},
              just={\TRule{\False}{\lor}[4]}, close]
          ]
        ]
        [\sFmla{\False}{\formula{A} \lor \formula{C}},
          just={\TRule{\False}{\land}[2]}, checked
          [\sFmla{\False}{\formula{A}},
            just={\TRule{\False}{\lor}[4]}, move by=2
            [\sFmla{\False}{\formula{C}},
              just={\TRule{\False}{\lor}[4]}, close]
          ]
        ]
      ]
      [\sFmla{\True}{\formula{B} \land \formula{C}},
        just={\TRule{\True}{\lor}[1]}
        [\sFmla{\False}{\formula{A} \lor \formula{B}},
          just={\TRule{\False}{\land}[2]}, checked
          [\sFmla{\False}{\formula{A}},
            just={\TRule{\False}{\lor}[4]}
            [\sFmla{\False}{\formula{B}},
              just={\TRule{\False}{\lor}[4]}
            ]
          ]
        ]
        [\sFmla{\False}{\formula{A} \lor \formula{C}},
          just={\TRule{\False}{\land}[2]}, checked
          [\sFmla{\False}{\formula{A}},
            just={\TRule{\False}{\lor}[4]}, move by=2
            [\sFmla{\False}{\formula{C}},
              just={\TRule{\False}{\lor}[4]}]
          ]
        ]
      ]
    ]
  ]
\end{oltableau}
નોંધો કે સ્પષ્ટતા માટે આપણે બીજી વખત $\TRule{\False}{\lor}$ લાગુ
કરવાનું પરિણામ નીચે ખસેડ્યું છે. આ કિસ્સામાં તેમ કરવું જરૂરી ન હોત,
કારણ કે સમર્થનો સરખાં જ હોત.

બે શાખાઓ હજી ખુલ્લી છે અને પંક્તિ~$3$ પરનું
$\sFmla{\True}{!B \land !C}$ ખરાના ચિહ્ન વિના છે. તેને
$\TRule{\True}{\land}$ લાગુ કરતાં બંધ !!{tableau} મળે છે:
\begin{oltableau}
  [\sFmla{\True}{\formula{A} \lor (\formula{B} \land \formula{C})},
    just=\TAss, checked
    [\sFmla{\False}{(\formula{A} \lor \formula{B}) \land
        (\formula{A} \lor \formula{C})}, just=\TAss, checked
      [\sFmla{\True}{\formula{A}}, just={\TRule{\True}{\lor}[1]}
        [\sFmla{\False}{\formula{A} \lor \formula{B}},
          just={\TRule{\False}{\land}[2]}, checked
          [\sFmla{\False}{\formula{A}},
            just={\TRule{\False}{\lor}[4]}
            [\sFmla{\False}{\formula{B}},
              just={\TRule{\False}{\lor}[4]}, close]
          ]
        ]
        [\sFmla{\False}{\formula{A} \lor \formula{C}},
          just={\TRule{\False}{\land}[2]}, checked
          [\sFmla{\False}{\formula{A}},
            just={\TRule{\False}{\lor}[4]}
            [\sFmla{\False}{\formula{C}},
              just={\TRule{\False}{\lor}[4]}, close]
          ]
        ]
      ]
      [\sFmla{\True}{\formula{B} \land \formula{C}},
        just={\TRule{\True}{\lor}[1]}, checked
        [\sFmla{\False}{\formula{A} \lor \formula{B}},
          just={\TRule{\False}{\land}[2]}, checked
          [\sFmla{\False}{\formula{A}},
            just={\TRule{\False}{\lor}[4]}
            [\sFmla{\False}{\formula{B}},
              just={\TRule{\False}{\lor}[4]}
              [\sFmla{\True}{\formula{B}},
                just={\TRule{\True}{\land}[3]}
                [\sFmla{\True}{\formula{C}},
                  just={\TRule{\True}{\land}[3]},close
                ]
              ]
            ]
          ]
        ]
        [\sFmla{\False}{\formula{A} \lor \formula{C}},
          just={\TRule{\False}{\land}[2]}, checked
          [\sFmla{\False}{\formula{A}},
            just={\TRule{\False}{\lor}[4]}
            [\sFmla{\False}{\formula{C}},
              just={\TRule{\False}{\lor}[4]}
              [\sFmla{\True}{\formula{B}},
                just={\TRule{\True}{\land}[3]}
                [\sFmla{\True}{\formula{C}},
                  just={\TRule{\True}{\land}[3]},close
                ]
              ]
            ]
          ]
        ]
      ]
    ]
  ]
\end{oltableau}

સરખામણી માટે, આ જ ધારણાઓના ગણ માટેનો એક બંધ !!{tableau} નીચે આપ્યો
છે, જેમાં નિયમો જુદા ક્રમે લાગુ કર્યા છે:
\begin{oltableau}
  [\sFmla{\True}{\formula{A} \lor (\formula{B} \land \formula{C})},
    just=\TAss, checked
    [\sFmla{\False}{(\formula{A} \lor \formula{B}) \land
        (\formula{A} \lor \formula{C})}, just=\TAss, checked
      [\sFmla{\False}{\formula{A} \lor \formula{B}},
        just={\TRule{\False}{\land}[2]}, checked
        [\sFmla{\False}{\formula{A}},
          just={\TRule{\False}{\lor}[3]}
          [\sFmla{\False}{\formula{B}},
            just={\TRule{\False}{\lor}[3]}
            [\sFmla{\True}{\formula{A}},
              just={\TRule{\True}{\lor}[1]},close
            ]
            [\sFmla{\True}{\formula{B} \land \formula{C}},
              just={\TRule{\True}{\lor}[1]}, checked
              [\sFmla{\True}{\formula{B}},
                just={\TRule{\True}{\land}[6]}
                [\sFmla{\True}{\formula{C}},
                  just={\TRule{\True}{\land}[6]},close
                ]
              ]
            ]
          ]
        ]
      ]
      [\sFmla{\False}{\formula{A} \lor \formula{C}},
        just={\TRule{\False}{\land}[2]}, checked
        [\sFmla{\False}{\formula{A}},
          just={\TRule{\False}{\lor}[3]}
          [\sFmla{\False}{\formula{C}},
            just={\TRule{\False}{\lor}[3]}
            [\sFmla{\True}{\formula{A}},
              just={\TRule{\True}{\lor}[1]},close
            ]
            [\sFmla{\True}{\formula{B} \land \formula{C}},
              just={\TRule{\True}{\lor}[1]}, checked
              [\sFmla{\True}{\formula{B}},
                just={\TRule{\True}{\land}[6]}
                [\sFmla{\True}{\formula{C}},
                  just={\TRule{\True}{\land}[6]},close
                ]
              ]
            ]
          ]
        ]
      ]
    ]
  ]
\end{oltableau}
\end{ex}

\begin{prob}
નીચેના માટે બંધ !!{tableau}s આપો:
\begin{enumerate}
\item $\sFmla{\True}{!A \land (!B \land !C)}, \sFmla{\False}{(!A \land !B) \land !C}$.
\item $\sFmla{\True}{!A \lor (!B \lor !C)}, \sFmla{\False}{(!A \lor !B) \lor !C}$.
\item $\sFmla{\True}{!A \lif (!B \lif !C)}, \sFmla{\False}{!B \lif (!A \lif !C)}$.
\item $\sFmla{\True}{!A}, \sFmla{\False}{\lnot\lnot !A}$.
\end{enumerate}
\end{prob}

\begin{prob}
નીચેના માટે બંધ !!{tableau}s આપો:
\begin{enumerate}
\item $\sFmla{\True}{(!A \lor !B) \lif !C}, \sFmla{\False}{!A \lif !C}$.
\item $\sFmla{\True}{(!A \lif !C) \land (!B \lif !C)}, \sFmla{\False}{(!A \lor !B) \lif !C}$.
\item $\sFmla{\False}{\lnot(!A \land \lnot !A)}$.
\item $\sFmla{\True}{!B \lif !A}, \sFmla{\False}{\lnot !A \lif \lnot !B}$.
\item $\sFmla{\False}{(!A \lif \lnot !A) \lif \lnot !A}$.
\item $\sFmla{\False}{\lnot(!A \lif !B) \lif \lnot !B}$.
\item $\sFmla{\True}{!A \lif !C}, \sFmla{\False}{\lnot (!A \land \lnot !C)}$.
\item $\sFmla{\True}{!A \land \lnot !C}, \sFmla{\False}{\lnot (!A \lif !C)}$.
\item $\sFmla{\True}{!A \lor !B}, \sFmla{\True}{\lnot !B},
  \sFmla{\False}{!A}$.
\item $\sFmla{\True}{\lnot !A \lor \lnot !B}, \sFmla{\False}{\lnot(!A \land !B)}$.
\item $\sFmla{\False}{(\lnot !A \land \lnot !B) \lif\lnot(!A \lor !B)}$.
\item $\sFmla{\False}{\lnot(!A \lor !B) \lif (\lnot !A \land \lnot !B)}$.
\end{enumerate}
\sourcecorrection{OLTAB-002}{મૂળની \texttt{proving-things.tex},
પંક્તિ~439માં અલ્પવિરામ અને $\lnot !B$ને પ્રથમ
$\sFmla{\True}{\cdot}$ની અંદર મૂક્યાં છે, જેથી તે એક સુનિર્મિત ચિહ્નિત
સૂત્ર રહેતું નથી. અભ્યાસના હેતુ મુજબ $\lnot !B$ને અલગ
$\sFmla{\True}{\lnot !B}$ ધારણા તરીકે પુનઃસ્થાપિત કર્યું છે.}
\end{prob}

\begin{prob}
નીચેના માટે બંધ !!{tableau}s આપો:
\begin{enumerate}
\item $\sFmla{\True}{\lnot(!A \lif !B)}, \sFmla{\False}{!A}$.
\item $\sFmla{\True}{\lnot(!A \land !B)}, \sFmla{\False}{\lnot !A \lor \lnot !B}$.
\item $\sFmla{\True}{!A \lif !B}, \sFmla{\False}{\lnot !A \lor !B}$.
\item $\sFmla{\False}{\lnot \lnot !A \lif !A}$.
\item $\sFmla{\True}{!A \lif !B}, \sFmla{\True}{\lnot !A \lif !B}, \sFmla{\False}{!B}$.
\item $\sFmla{\True}{(!A \land !B) \lif !C}, \sFmla{\False}{(!A \lif !C) \lor (!B \lif !C)}$.
\item $\sFmla{\True}{(!A \lif !B) \lif !A}, \sFmla{\False}{!A}$.
\item $\sFmla{\False}{(!A \lif !B) \lor (!B \lif !C)}$.
\end{enumerate}
\end{prob}

\end{document}
